to the series @@M0@@, where @@M1@@ is the class of @@M2@@ in @@M3@@. Its kernel is the group @@M4@@ of formal series such that:
@@M5@@
we therefore have an isomorphism, functorial in @@M6@@:
@@M7@@

\sgaStableTarget{sga6.e05.lem.1-12}\textbf{Lemma 1.12.} \uline{Let} @@M8@@ \uline{be a functor on} @@M9@@; \uline{then the canonical homomorphism}:
@@M10@@
\uline{is an isomorphism}.

The homomorphism @@M11@@ is defined as the composite:
@@M12@@
the last arrow being defined by the homomorphism @@M13@@

We shall construct a morphism in the opposite direction:
@@M14@@
For this, we shall show that every morphism of functors from @@M15@@ to @@M16@@ factors through the canonical morphism @@M17@@, which, by 1.11, is equivalent to showing that it vanishes on @@M18@@.

Let @@M19@@, @@M20@@; we define a homomorphism of
